% Generated draft: requires mathematical and source review.
\documentclass[11pt]{article}
\usepackage[margin=25mm]{geometry}
\usepackage{fontspec}
\setmainfont{Latin Modern Roman}
\usepackage{amsmath,amssymb,mathtools,unicode-math}
\setmathfont{Latin Modern Math}
\usepackage{enumitem,needspace,xurl,hyperref}
\setlength{\parindent}{0pt}
\setlength{\parskip}{5pt}
\setlength{\emergencystretch}{3em}
\newcommand{\N}{\mathbb{N}}
\newcommand{\Z}{\mathbb{Z}}
\newcommand{\Q}{\mathbb{Q}}
\newcommand{\R}{\mathbb{R}}
\newcommand{\C}{\mathbb{C}}
\newcommand{\F}{\mathbb{F}}
\newcommand{\A}{\mathbb{A}}
\newcommand{\PP}{\mathbb{P}}
\newcommand{\E}{\mathbb{E}}
\newcommand{\eps}{\varepsilon}
\newcommand{\dd}{\,\mathrm{d}}
\begin{document}
\section*{UnsolvedMath ID: 6700072}
\textbf{Source problem code:} catalogue entry
\subsection{Definitions and mathematical statement}
The main conjecture, ``Stability of Periodic Slabs,'' says that the only $\Z^{n-3}$-invariant mean-convex domains in $\R^n$ with disconnected boundaries are slabs between parallel hyperplanes. In the usual translation-periodic reading, $n\ge3$ and a rank-$(n-3)$ lattice of translations preserves a domain $Y$; the proposed conclusion is
\[
 Y=\{x\in\R^n:a<\langle x,v\rangle<b\}
\]
for a unit vector $v$ and real $a<b$, with the appropriate boundary convention. The source does not specify the periodic action, boundary regularity, or sign and non-strict/strict meaning of mean convexity. Translation periodicity is an explicit conventional reading, not a supplied definition.

The passage then discusses a different class of ``thick'' mean-convex-domain constructions. It considers a complete Riemannian $n$-manifold $X$ satisfying three conditions:
\begin{enumerate}
\item $X$ is connected at infinity; the precise end convention is not stated.
\item Its boundary at infinity has infinite $(n-1)$-volume in the source's stated sense: every proper continuous $f:X\to\R_+$ satisfies
\[
 \limsup_{t\to\infty}\operatorname{vol}_{n-1}(f^{-1}(t))=\infty.
\]
The measure of level sets of arbitrary continuous functions needs an explicit convention.
\item $X$ is locally ``$(n-2)$-thick'': there exist $\varepsilon>0$, $\alpha>1$, and $c>0$ such that every $(n-2)$-cycle $B$ with
\[
 \operatorname{diam}(B)\le\varepsilon,\qquad
 \operatorname{vol}_{n-2}(B)\le\varepsilon^{n-2}
\]
bounds an $(n-1)$-chain $C$ satisfying
\[
 \partial C=B,\qquad
 \operatorname{vol}_{n-1}(C)\le
                  c\bigl(\operatorname{vol}_{n-2}(B)\bigr)^\alpha.
\]
Cycle coefficients, chain category, and volume definitions are not supplied.
\end{enumerate}
Under these three conditions the source asserts that every compact $Y_0\subset X$ is contained in a smooth compact mean-convex domain $Y_1\subset X$. It contrasts these thick domains with thin/narrow ones surrounding codimension-two subsets. This construction statement is retained as adjacent context, not a proof of the periodic-slab conjecture.

The source says that manifolds uniformly bi-Lipschitz homeomorphic to $\R^n$, or to ``any Riemannian homogeneous space,'' satisfy all three conditions. The latter blanket phrase needs qualifications: completeness and homogeneity alone do not specify noncompactness or connectedness at infinity, and compact homogeneous spaces do not have the stated unbounded-level-set property. No intended restriction is guessed. The extraction ends at ``Sketch of the,'' so the asserted construction's proof is absent.
\subsection{Short English statement}
Must a periodic mean-convex Euclidean domain with disconnected boundary be a slab? The source also includes a separate compact mean-convex-enclosure assertion under end, level-set-area, and local filling assumptions. Both statements are retained with their distinct hypotheses.

The supplied summary reports no resolution of periodic-slab rigidity and gives only general mean-curvature/topology context. The OPEN-TRIAGE label and 2017 passage are retained with current status marked for review. The adjacent homogeneous-space claim and truncated proof do not independently verify either statement.
\subsection{Sources}
\begin{itemize}
\item Local UnsolvedMath registry, permanent ID 6700072, source code catalogue entry, scalar-curvature Question [?76], ``Stability of Periodic Slabs,'' proposed by Misha Gromov, 2017. Dataset curation and metadata: Ulam AI and UnsolvedMath Contributors, CC BY 4.0.
\item Misha Gromov, \emph{101 Questions, Problems and Conjectures around Scalar Curvature} (2017), printed tag [?76], occurrence 75, PDF page 76, including the extracted thick-domain passage, \url{https://www.ihes.fr/~gromov/wp-content/uploads/2018/08/101-problemsOct1-2017.pdf}; normalized pdf-text extraction, accessed 2026-07-29.
\item No exact later resolution or proof of the adjacent construction is supplied. Source-specific provenance: the author-hosted PDF is incomplete and publicly available; current status and statement-reuse terms NEEDS\_REVIEW, separately from the dataset licence.
\end{itemize}
\end{document}
